% arXiv change notes
% 29/08/24 Minor clarification of r(z) function

\documentclass[a4paper,10pt]{amsart}
%\documentclass[referee,sn-basic]{sn-jnl}% referee option is meant for double line spacing
\usepackage{fourier}
\usepackage{biblatex}

\addbibresource{Phyllotaxis.bib}

\usepackage{graphicx}%

\usepackage{amsmath,amssymb,amsfonts}%
\usepackage[foot]{amsaddr}

\newlength{\jfigwidth}
\setlength{\jfigwidth}{\textwidth}
\newcommand{\jpgfig}[3]{\jdofig{#1}{#2}{#3}{Figures}{.jpg}}
\newcommand{\pdffig}[3]{\jdofig{#1}{#2}{#3}{Figures}{.pdf}}
\newcommand{\jdofig}[5]{
	\begin{figure}\centering\includegraphics[width=#3\jfigwidth]{#4/#1#5} \caption{#2}\label{fig:#1}\end{figure}
}

\usepackage[colorinlistoftodos, textsize=tiny]{todonotes}
\begin{document}

\title[Critical values]{Critical values for preserving Fibonacci structure in disk-stacking models} 
\author{Jonathan Swinton}
\address{ORCID 0000-0002-3196-9653}
\email{jonathan@swintons.net}


\listoftodos % This generates the summary of all notes
\newpage

\begin{abstract}
This is the outline of a paper exploring the computational and hopefully analytic properties of the threshold at which disk-stacking models fail to preserve Fibonacci structure.

\end{abstract}
\maketitle 

%\listoffixmes
 \section{Introduction}
This note explores the dynamics of Schwendener disk-stacking models reintroduced into phyllotaxis by Douady and Golé. A tutorial introduction to these models, along with further references can be found in my textbook~\cite{swintonMathematicalPhyllotaxis2025}. They are of particular interest, in part because they connect with the known biology of plant organ formation, but more relevantly here because they are capable of showing Fibonacci structure. In particular, they have a geometric parameter which is the ratio of the size of the disks being stacked to the arena in which they are being stacked in, and when this parameter is smoothly and slowly changed, the resulting arrangement of disks exhibits families of spirals characterised by increasing pairs of Fibonacci numbers. If the geometric parameter is changed too quickly, this property is lost, and in~\cite{swintonDiskstackingModelsAre2024} I explored 
  with some success whether this, in conjunction with the introduction of noise, could explain observed biological data which often includes non-Fibonacci spiral counts and sometimes the complete loss of spiral structure.
 That paper reconfirmed the observations of Douady and Golé that there is a non-dimensional threshold for the rate of change of the geometric parameter above which non-Fibonacci transitions occur, but at present this threshold is poorly characterised numerically and not at all analytically. 
 
  
 
 
\section{Methods}
Instead of the fixed-circumference cylinder bur  changing stacking-disk radius of previous models, here I keep the radius of each stacking disk constant but vary the circumference of the periodic arena in which it is stacked. In the cylindrical model, the disk radius was modelled as $r(z)$, where $z$ was vertical height on the cylinder, and $r(z)$ was in simple cases linear in $z$. Then $r'=dr/dz$ was a non-dimensional measure of the rate of change of disk size. Here, instead we keep $r(z)==1$ but vary the cone circumference linearly with $z$. This has numerical advantages and is perhaps easier to interpret visually, and $r'$ becomes roughly $1/\tan \theta$ where $\theta$ is the semi-angle of the cone.\todo{is this right?}


 
 The simulation code relies on a Wolfram notebook at \url{https://github.com/js229/GeometricalPhyllotaxis/blob/master/Mathematica/Packages/CoinStacking.nb}, and the simulations themselves are run by \url{https://github.com/js229/GeometricalPhyllotaxis/blob/master/LaTeX/Critical%20values/CriticalR.nb}; the first line of the latter notebook needs to point to the  location of the first.
 
 Figure~\ref{fig:jCRfibonacci} gives an example of the simulation. 
 \pdffig{jCRfibonacci}{A stacked-disk simulation with semi-cone angle $0.12\pi$ radians. The disk-stacking takes place on  a cone, with periodic copies of disks indicated by the open circles, and at each iteration the next disk is placed at the closest possible position to the cone apex, without overlapping any previously placed disks. 
 	Each disk is joined to the two disks which support it from below by a red line to one on its right and a blue line to the one on its left; disks not joined by lines are not touching. 
 	 This pattern has a $(13,21)$ parastichy at the top rim, as measured by counting the 13   blue lines running upwards and to the left, taking account of the cone periodicity, and the 21 red lines running upwars and to the right.
 	 }{1}%  
 
 Parastichy lines are constructed from the finished simulation as explained in the Figure caption, and parastichy counts are calculated as how many different parastichy lines can be seen in at the top rim of the cone where the simulation has been terminated.
 
 \section{Results}
 \pdffig{jCRTransition}{First results: parastichy counts as a function of the cone semi-angle. In each case the simulation was run to fill a cone with upper rim circumference $2\pi \cdot 55$}{1}%  
 Figure~\ref{fig:jCRTransition} shows the parastichy counts as a function of the speed at which geometry changes as measured by the semi-angle of the cone. For narrow cones with semi-angle of about $0.155\pi$, each transition follows the Fibonacci rule associated with a pattern locally near a hexagonal packing as discussed in~\cite{swintonMathematicalPhyllotaxis2025}. 
 For values of the cone semi-angle  beyond around $0.155\pi$, the reliable Fibonacci transitions break down, although there is evidence of continued structure until about $0.165\pi$.
 
 The region near $0.166$ where the cone semiangle is $~\pi/6$ shows a coherent structure (Figure~\ref{fig:jCRIntermediate}). On examination of the resulting disk pattern this nay because the expanding cone is very nearly able to accomodate a repeated lattice-like pattern. 
 
 I would expect the details of the parastichy counts and disk patterns above transition to be quite dependent on the point at which the pattern is evaluated.
 
 Figure~\ref{fig:jCREnd} gives an example well past the transition. 
 
\pdffig{jCRIntermediate}{Resulting pattern is close to hexagonal when the cone semi-angle is close to $\pi/6$. }{1}%  
\pdffig{jCREnd}{Resulting pattern when the cone semi-angle is close to $\pi/5.5$. The vertical parastichy lines reflect the stacked column structure which  Doaudy and Golé call quasi-symmetric.
}{1}%  


 
\clearpage
\section{Analytic approach}
Suppose that we took the chain at each height as coming from the van Iterson touching-circle lattice from the corresponding geometry ratio?
\section{Outstanding issues}
\begin{itemize}

	\item Is the threshold cone angle the same size even starting from large Fibonacci patterns
	\item Make the cone piecewise linear and explore how to reliably generate eg Lucas numbers
	\item Do the results hold for longer runs
	\item Add noise along the lines of the previous paper
	\item fine structure of transition - expect to see lots of islands of structure when we look carefully
	\item Relate to cylinder paper results and at a minimum connect $\theta$ to $r'$.
	\item Explore geometry at transition to derive a condition for the critical $\theta$.
	\item 
\end{itemize}


\section{Some other topics}
Out of scope for this paper, but of interest
\begin{itemize}
	\item  log term dynamics on a cone - Figure~\ref{fig:jCRLongTerm} looks explicable.
\item	 long term dynamics on a cylinder, ie with no change in relative geometry. This also has lots of good open questions such as whether all attractors are rhombic tilings.

\end{itemize}
\pdffig{jCRLongTerm}{Long term Fibonacci transitions have intriguing structure. Diagram as before, but without disks and parastichy lines recoloured to emphasis topological structure. 
}{1}%  


\section*{Statements and Declarations} 
I did not receive support from any organization for this work and have no relevant financial interests to disclose. 

\printbibliography
\section{Document history}
Last compiled\ \today.
 Mathematica code used to generate the Figures, along with a version-controlled copy of the source text, can be found at \url{https://github.com/js229/GeometricalPhyllotaxis/tree/master/LaTeX/Critical%20values}.
\end{document}



